\section{The exact ind-path formula and the manuscript corrections}
\label{sec:path-manuscript}

The conventions on universes, arbitrary interior vertical arrows and
endpoint identities are those of C.1. Write
\[
 H_{a,b}(X)=\operatorname{Hom}_{M^X[\Sigma_X^{-1}]}(a_X,b_X),
 \qquad
 \Sigma_X=\{\eta\mid \eta_x\in\Sigma\text{ for every }x\in X\}.
\]
The category $M^X$ and its marking are $U$-small for every
$U$-small $X$. In particular the quantifier on $X$ below includes
infinite and empty categories.

\begin{claim}{C.20}{Proved}{D.K5 with two complete proofs and its exact marking convention}
Let $t:\mathbb N^{\mathrm{op}}\to\mathrm{Ty}$ be a sequence
such that every type receives a map from some $t_n$, and put
$K_n=F_{t_n}(a,b)$ with its actual pullback transitions. For every
small $M$, wide composition-closed $\Sigma$, and small $X$, effect
of a localization word induces a natural bijection
\[
 \operatorname{colim}_{n}\pi_0\operatorname{Fun}(X,K_n)
       \xrightarrow{\ \sim\ }H_{a,b}(X).                 \tag{C.20.1}
\]
It respects concatenation, the constant paths, and restriction in
$X$. Containing all identities and being closed under composition
are exactly the conditions on a class of arrows that make these
path categories stable under every type pullback.
\end{claim}
\begin{proof}[First proof: words, ladders, and actual transitions]
A path of type $\tau$ in $M^X$ from $a_X$ to $b_X$ is literally a
functor $X\to F_\tau(a,b)$. The exponential law identifies both
its objects and its transformations: it is a functor
$X\times I_\tau\to M$, negative edges are pointwise marked,
and the endpoint diagrams and the endpoint components are fixed.
Thus
\[
 F_\tau^{(M^X,\Sigma_X)}(a_X,b_X)
       =\operatorname{Fun}(X,F_\tau^{(M,\Sigma)}(a,b))       \tag{C.20.2}
\]
up to the canonical exponential-law isomorphism, commuting with
every type map.

We check the component theorem used in D.K5. Compose positive
edges and inverses of marked negative edges in the ordinary
localization. A type pullback preserves this effect. Every ladder
preserves it as well: a positive square gives the usual composite
equation; a negative square gives the same equation after
cancelling its two invertible horizontal arrows. No vertical arrow
is inverted. Conversely form a category with objects those of $M^X$
and with morphisms the components of its all-type path categories.
Literal concatenation defines composition. Two successive forward
edges and their composite are joined by the type map selecting the
two outer vertices. Inserting an identity is joined by collapse.
For a marked $w$, the two backtracks are joined to the constant
path by the ladder with middle component $w$ and endpoint
components identities, in the directions written in C.7. Hence
the structural functor inverts $\Sigma_X$. For any ordinary functor
out of $M^X$ inverting $\Sigma_X$, evaluation of the same words
defines a unique factorization, by the square calculation above.
This proves the universal property of localization and the exact
component theorem, including its specified effect map.

Each localization word has one finite type, even when $X$ is
infinite. Pull it back along $t_n\to\tau$ to obtain surjectivity
in (C.20.1). To check injectivity, an equality in the component
category is witnessed by a finite zigzag of total-path arrows.
Collect the finitely many types involved. Their concatenation maps
to each of them, and some $t_N$ maps to that concatenation. Pull
each ladder back to $t_N$. At a type-change arrow
$(u,\eta):c\to d u$, the two resulting pullbacks of $d$ may use
different maps $t_N\to\operatorname{type}(d)$. These maps lie in
the same nonempty $P(t_N,\operatorname{type}(d))$. C.2 connects
them by a finite zigzag of pointwise transformations; its strict
action, followed by $\operatorname{Fun}(X,-)$, connects the
corresponding diagrams. At the two ends the chosen maps from
$t_N$ may likewise differ from the actual transition maps of $t$.
The same connectedness supplies those two comparisons. Enlarge
$N$ first to dominate the original stages; all required maps then
exist. We have obtained equality at an actual common stage of the
given sequence, which proves injectivity. Every construction used
effect, restriction, or concatenation, so the stated naturalities
and the units follow from their explicit formulas.
\end{proof}

\begin{proof}[Second proof: the specified homotopy-colimit span]
Apply C.3 to $(M^X,\Sigma_X)$, which meets the same hypotheses.
Its specified span is a pair of weak equivalences, so taking
components identifies the all-type component category with the
components of its filtered categorical colimit. The functor
$\pi_0:\mathrm{Cat}\to\mathrm{Set}$ is left adjoint to the
discrete-category functor and therefore preserves that colimit.
Use (C.20.2) and the component universal property just checked.
The composite sends every word to its effect, so it is precisely
(C.20.1), not another bijection.

Finally the pullback of a negative one-edge path along
$-\to\varnothing$ requires every identity to be marked; pullback
along $-\to--$ requires the composite of every composable marked
pair to be marked. These two requirements suffice, because the
image of a negative edge is a negative string or an identity.
For a class $S$ without these closure properties, its generated
wide subcategory $\langle S\rangle$ has the same ordinary and
$\infty$-categorical localization: a functor inverts $S$ exactly
when it inverts all identities and composites generated by $S$.
The two universal properties supply inverse comparison functors
under $M$. The formula then uses $\langle S\rangle$, with those
specified comparison functors. Without closing the marking the
type diagram itself need not exist. For example in $[2]$, marking
the identities and $0\to1,1\to2$ but omitting $0\to2$ makes
pullback along $-\to--$ fail. With $M=*$ and no marked identity,
every mixed-type fibre is empty while the localized Hom set is a
singleton. These finite failures are checked in the accompanying
program.
\end{proof}

\begin{claim}{C.21}{Proved}{Complete classification of projective sequences of types}
For a projective sequence $t$, let $r_n$ be the number of maximal
constant-sign runs of $t_n$, with zero for the empty word. The
following three conditions are equivalent: the $r_n$ are unbounded;
every type receives a map from some $t_n$; and (C.20.1) is a
bijection for every small relative category, all endpoint pairs,
and every small $X$.

Every sequence outside this class is eventually, after passing to
$\mathrm{Cat}_{\mathrm{hot}}$ on its path fibres, constant at the
fibre of one finite alternating word $\sigma$. More precisely,
for every $(M,\Sigma),a,b,X$ there is a canonical bijection
\[
 \operatorname{colim}_n\pi_0\operatorname{Fun}(X,F_{t_n}(a,b))
      \cong\pi_0\operatorname{Fun}(X,F_\sigma(a,b)),        \tag{C.21.1}
\]
commuting with effect in the localization. Consequently, for any
fixed data outside the universal domain, (C.20.1) holds exactly
when every localization word from $a_X$ to $b_X$ has a
$\sigma$-representative and any two such representatives with
the same effect are connected by endpoint-fixed $\sigma$-ladders.
These are respectively the exact existence and equality conditions
for this finite type; no claim that they always hold is made.
\end{claim}
\begin{proof}
The letter-map description in C.11 shows that $\tau\to\rho$
implies $r(\rho)\leq r(\tau)$: successive runs of opposite sign
must occupy successive, distinct runs of the source letters.
Thus $r_n$ is nondecreasing. Arbitrarily long alternating run
strings contain the letters of each fixed target word in order,
allowing repetitions within a sign run. The resulting letter map
is the required type map. Conversely maps to all alternating words
force unbounded run counts. C.20 proves sufficiency for the formula.

For necessity, suppose $r_n\leq r$. Let $M$ be the alternating
fence poset with vertices $0,\ldots,r+1$, mark all its arrows,
and take endpoints $0,r+1$ and $X=*$. Its nerve is a subdivided
interval, so its groupoid completion has exactly one arrow between
these endpoints. A same-sign run of a path in this fence can
contain at most one nonidentity arrow: the fence has no composable
pair of nonidentity arrows in either direction. It therefore moves
at most one unit in undirected distance. No path with at most $r$
runs joins vertices at distance $r+1$. Every $F_{t_n}(0,r+1)$
is empty, so the left side of (C.20.1) is empty and the right side
is a singleton. This is a finite counterexample for every finite
bound, including $r=0$.

For the stronger individual statement, bounded monotone $r_n$
eventually equals $r$. At equal positive run count the first signs
of two words admitting a type map must agree; otherwise embedding
the alternating target-run pattern would require an additional
source run. The reduced alternating word $\sigma$ is therefore
constant on this tail. There are type maps in both directions
between $\sigma$ and each $t_n$, obtained by selecting or
repeating the letters in each run. Their composites and the
identities belong to the same connected posets $P(t_n,t_n)$ and
$P(\sigma,\sigma)$. C.2 shows that the induced functors on path
fibres are inverse in $\mathrm{Cat}_{\mathrm{hot}}$. Different
choices give the same morphisms there. For each transition the
two maps $t_m\to\sigma$ are again connected in their mapping
poset. Hence these equivalences identify the whole tail diagram
in $\mathrm{Cat}_{\mathrm{hot}}$ with the constant one at
$F_\sigma$. Apply $\pi_0\operatorname{Fun}(X,-)$ to obtain
(C.21.1), including $r=0$. The two conditions in the statement are
exactly surjectivity and injectivity of its specified effect map,
as follows from the localization presentation proved in C.20.
\end{proof}

\begin{claim}{C.22}{Proved}{The type formalism and its actual colimits}
The corrected type category at VII ms.~26--32 has
endpoint-preserving, nondecreasing interval functors, and its path
construction is the fibred construction of C.1. Its concatenation
is strictly associative and unital. Its empty word is terminal.
The opposite type category is neither filtered nor totally
$W_\infty$-aspherical. In contrast every sequential path colimit
$C_\infty(t)$ is indexed covariantly by the filtered category
$\mathbb N$, regardless of cofinality of $t$.

For a sequence of full subcategory inclusions $K_n\subset K_{n+1}$,
every $X$ with finitely many objects satisfies
\[
 \operatorname{colim}_n\operatorname{Fun}(X,K_n)
       \cong\operatorname{Fun}(X,\operatorname{colim}_nK_n).
                                                               \tag{C.22.1}
\]
For arbitrary filtered diagrams the same assertion holds for
categories $X$ with finitely many arrows. These are statements
in $\mathrm{Cat}$ and on taking components; the map to derived
mapping classes is the further specified comparison $c_X$ of C2.
\end{claim}
\begin{proof}
Endpoint preservation is exactly what makes endpoint evaluation a
functor on the total category: the endpoint component of
$c\Rightarrow d u$ has target $d(0)$ or $d(n)$ precisely when $u$
fixes the endpoints. The monotonicity restriction is the definition
adopted at ms.~31. It makes each source edge map to a same-sign
string or an identity and gives the letter correspondence of C.11.
The functor to $\mathrm{Cat}$ is faithful, but the inclusion into
all endpoint-preserving interval functors need not be full. For
example the vertex sequence $(0,1,0,1,2,3)$ defines an
endpoint-preserving functor $I_{+-+-+}\to I_{+-+}$: its five
arrow images are respectively $0\to1$, $0\to1$, $0\to1$,
$2\to1$, and $2\to3$. It decreases between its second and
third vertices and therefore is not a map of types.

For ms.~48--58 the pushout of two interval categories along their
common endpoint is their concatenated interval: specifying a
functor out of the pushout specifies exactly two functors with
the same value at the junction. The same assertion for natural
transformations gives the stated pullback of path categories.
Blockwise gluing of vertex maps obeys
$(gf)*(g'f')=(g*g')(f*f')$, with strict units and associativity.
An arrow of types is uniquely the concatenation of its restrictions
to the individual source edges. Each restriction is
$\varepsilon\to\varepsilon^{l}$ for a unique $l\geq0$.
This proves the free-monoid description of arrows at ms.~50'--51.
For a composable pair, refine each first-edge image into the
successive second-edge image lengths. This yields exactly the
functions $[1,l]\to\mathbb N$ at ms.~52--53. Iterating this
description proves the analogous description for every nerve
degree at ms.~54--58, including $l=0$; composition adds the image
lengths, and identity insertion assigns every length one. This
verifies the combinatorial assertions actually used in the path
construction.

There is one map from every interval to the empty type. The
parallel maps of C.11 admit no common equalizing completion in
$\mathrm{Ty}^{\mathrm{op}}$, so that category is not filtered;
C.11 gives its stronger failure of total asphericity. An actual
projective sequence, however, acts contravariantly on paths and
therefore gives a covariant $\mathbb N$-diagram. Its colimit is a
homotopy colimit by the proof in C.3. The transition at VII
ms.~102 must accordingly be $C_{t_{n-1}}\to C_{t_n}$.

For (C.22.1) the images of the finitely many objects of $X$ lie in
one stage. Fullness puts all images of arrows there, even when
there are infinitely many arrows, and faithfulness retains all
their equations. The components of a transformation form another
finite list; the same argument retains them and all their
naturality equations. This proves the category isomorphism. For
arbitrary filtered diagrams, a finite-arrow category has finitely
many generators and finitely many composition equations, so the
same data-and-equations argument applies. A zigzag connecting two
objects is finite; hence the component assertions follow too.
\end{proof}

\begin{claim}{C.23}{Proved}{VII \S10: the uniform criterion and its path-specific replacement}
For small $u:J\to T$, the projection of fibred total categories
\[
 p_F:\int_J^{\mathrm{fib}}Fu^{\mathrm{op}}
               \longrightarrow\int_T^{\mathrm{fib}}F
\]
is a weak equivalence for every strict $F:T^{\mathrm{op}}\to
\mathrm{Cat}$ exactly when $B(u/t)$ is weakly contractible for
every $t\in T$. The parallel-pullback homotopy, preorder filtration
and countability hypotheses of VII ms.~99--101 do not imply this
comparison. For path fibres over $\mathrm{Ty}$ the specified
comparison does hold, by C.2--C.3. The projection
$\mathrm{Ty}\to R$ is neither smooth nor proper, even for the
connected-component localizer, so neither sufficient criterion
proposed at VII ms.~112 supplies that path comparison.
\end{claim}
\begin{proof}
For $c=(t,z)$ in the target of $p_F$, forget fibre data in $p_F/c$.
This gives a functor to $u/t$. An object over
$(j,\alpha:u(j)\to t)$ is an arrow $x\to F(\alpha)z$ in
$F(u(j))$. Selecting its terminal arrow $1_{F(\alpha)z}$ is a
functorial section, and the given arrow supplies a natural
transformation from the identity to section-after-projection.
The assertion about the section on arrows is the equation
$F(\alpha)=F(u(v))F(\alpha')$ for a comma arrow
$v:(j,\alpha)\to(j',\alpha')$. Thus
$B(p_F/c)\simeq B(u/t)$. Quillen's comma criterion proves the
forward implication from the displayed condition. Conversely take
the discrete presheaf $F(s)=\operatorname{Hom}_T(s,t)$.
Its total category is $T/t$, with terminal object, and its pulled
back total category is $u/t$. This proves necessity with the
universal quantifier over $F$, precisely as in PATHS A.3.

For the failure take $T=\mathbf B C_2$, constant terminal fibre,
and $u:\mathbb N^{\mathrm{op}}\to T$ sending every arrow to
the identity. All parallel pullbacks coincide, the reachability
preorder is a singleton, and the countability condition holds.
Nevertheless the total category is $\mathbf B C_2$, the pulled
back total category is $\mathbb N^{\mathrm{op}}$, and its
filtered colimit is a point. The comparison induces
$0\to C_2$ on fundamental groups. This confirms both PATHS A.4
and Grothendieck's own example at ms.~109.

For failure of smoothness and properness let $\sigma=+-+$ and
$\tau=+$. Both have only the identity endomorphism in
$\mathrm{Ty}$. There are exactly two type maps
$\sigma\to\tau$, with vertex tuples $(0,0,0,1)$ and
$(0,1,1,1)$. Above the single base arrow $p\sigma\to p\tau$
in $R$, its incoming-lift category at $\tau$ has these two
objects; a morphism would be an endomorphism of $\sigma$, so
only identity arrows occur. Its outgoing-lift category at $\sigma$
is likewise the discrete two-object category, using endomorphisms
of $\tau$. In the criterion of
\cite[Theorem~1.2(b)]{CpathsMaltsiniotis}, the first category is
the comma category of the fibre inclusion over
$(\tau,p\sigma\to p\tau)$; the second is its opposite-category
dual. Both fail even connectedness. This proves the asserted
failure with the exact variance and gives the calculation behind
the manuscript's statement. C.2 supplies instead a contractible
\emph{poset of transformations between} these type maps; its
strict action and the coherent descent in C.3 give the path
comparison despite these two discrete lifting categories.
\end{proof}

\begin{claim}{C.24}{Refuted, with precise replacement}{V ms.~46--49: fibrant unions and ind-isomorphisms}
Filtered colimits of Thomason-fibrant categories are
Thomason-fibrant, and objectwise weak equivalences of actual
filtered diagrams induce weak equivalences of colimits. There is,
however, a sequence of finite local Thomason-fibrant categories
with full inclusion transitions whose union $T$ is not local.
The union has an endofunctor $s$ inducing the identity ind-morphism
in $\operatorname{Ind}(\mathrm{Cat}_{\mathrm{hot}})$ but not a
weak equivalence. Thus the proposed implications at V ms.~47 and
the uniqueness claim at ms.~48--49 fail even with fibrant stages
and a fibrant union.
\end{claim}
\begin{proof}
For fibrancy, each generating Thomason horn
$c\operatorname{Sd}^2\Lambda_i^r\to
c\operatorname{Sd}^2\Delta^r$ is a map between finite categories.
A lifting problem into a filtered colimit is represented at one
stage, where a filler exists. Filtered colimits commute with nerves
and preserve simplicial weak equivalences, by the finite-data Kan
replacement argument in C.3. These prove the two positive claims.
An actual sequence admits compatible fibrant replacements:
choose trivial cofibrations $K_n\to R_n$ with fibrant targets,
then lift against $R_{n+1}\to *$ to extend
$K_n\to K_{n+1}\to R_{n+1}$ over $R_n$.
The freely generated sequential arrows need no further coherence
equations. This verifies the replacement construction at V ms.~43
without asserting any additional local mapping property.

Here is an independent check of CATHOT E.4. Let
$G_n=\operatorname{Sym}\{0,\ldots,n\}$ and let $T_n$ have
objects $0,\ldots,n$ and
\[
 \operatorname{Hom}_{T_n}(i,j)=
 \begin{cases}G_j&i\leq j,\\\varnothing&i>j.\end{cases}
 \qquad (h:j\to k)(g:i\to j)=h\,\iota_{jk}(g).
\]
The inclusions of permutation groups make this associative.
The identity labels give identities. Selecting object $n$ defines
$j_n:\mathbf B G_n\to T_n$, and inclusion of arrow labels
defines $r_n:T_n\to\mathbf B G_n$. Then $r_nj_n=1$ and the
identity labels give $1\Rightarrow j_nr_n$. Thus each $T_n$
is local by its specified categorical homotopy equivalence to a
groupoid. The inclusions $T_n\subset T_{n+1}$ are full.

Each $T_n$ is also Thomason-fibrant. Given a generating horn
$A\subset B$ and $A\to T_n$, extend its composite to
$\mathbf B G_n$ across $B$, which is possible since a groupoid
has Kan nerve and hence a fibrant Thomason nerve. Keep the old
objects and send every new object to $n$; on each new arrow use
the label of the extended functor. The horn inclusion is a full
sieve: twice subdividing a simplicial subcomplex gives the
inclusion of its downward-closed face-chain poset. Thus no arrow
goes from a new object to an old one. Every remaining new arrow
has an allowed target, and the label multiplication verifies its
mixed composition equations. This is a strict filler. The union
$T$ is therefore fibrant by the first paragraph.

Define $s:T\to T$ by $i\mapsto i+1$ and by shifting every
permutation one place, fixing zero. It respects the displayed
composition formula. For fixed $n$, put $m=n+1$. The inclusion
and shifted functors $T_n\to T_m$ each have an identity-label
transformation to the functor constant on objects at $m$ with
the corresponding arrow labels. The latter two functors are
naturally isomorphic by $c=(0\ 1\ \cdots\ n+1)$, since
$cgc^{-1}=s(g)$ for every $g\in G_n$. This gives a three-arrow
zigzag from inclusion to shift in
$\operatorname{Fun}(T_n,T_{n+1})$. Hence $s$ induces the
identity ind-morphism.

The $r_n$ are strictly compatible with the group inclusions and
give a weak equivalence $BT\to B G$, where
$G=\operatorname{Sym}_{\mathrm{fin}}(\mathbb N)$, by the first
paragraph. Under it $Bs$ induces the group shift. Its image fixes
zero and omits the transposition $(0\ 1)$, so it is not a
fundamental-group isomorphism. Consequently $Bs$ is not a weak
equivalence and cannot be homotopic to the identity. The union is
not local: a functor $\mathbf B G\to T$ chooses one object $i$
and has its entire endomorphism image in the finite group $G_i$.
Its classifying map cannot be an inverse to $BT\simeq BG$,
even up to unpointed homotopy, since conjugation preserves finiteness
of image. The local comparison for that source therefore fails.
The attached program independently checks all category equations,
the retractions and the three transformations through $n=3$.

The unbounded-zigzag-length mechanism mentioned at ms.~49 also
occurs concretely. Finite intervals in the infinite alternating
fence contract by successively deleting an endpoint, giving finite
local full stages. A categorical homotopy of length $l$ from its
union's identity to a constant functor would give an undirected
path of length at most $l$ from every integer vertex to one fixed
vertex. The fence has unbounded distances, so no such homotopy
exists. Finally fibrancy alone does not replace a zigzag by one
arrow: the terminal fork $0\to1\leftarrow2$ is fibrant, since
every functor from a full sieve extends by assigning its terminal
object to new objects. Its objects $0$ and $2$ are joined by a
two-arrow zigzag and have no arrow in either direction.
\end{proof}

\begin{claim}{C.25}{Refuted, with an explicit finite obstruction}{VII ms.~129--132: invariance is additional structure}
Weak-equivalence invariance of $H_{a,b}$ can fail between finite
test categories for a finite relative category with every arrow
marked. Thus neither the existence of the mapping space in C.7
nor weak equivalence of a path category to a Thomason-fibrant
category implies this invariance. The canonical mapping-space
interpretation itself always holds, before making any invariance
assumption on $H_{a,b}$.
\end{claim}
\begin{proof}
We verify PATHS D.5 completely. Let $Y$ be the four-object crown
with minima $r_0,r_1$, maxima $q_0,q_1$, and all four arrows
$r_i\to q_j$. Adjoin two incomparable objects $s,n$, each above
every object of $Y$, to obtain a six-object poset $M$. Mark all
arrows. Their nerves are respectively the four-cycle and the
join of three two-point sets, hence $BY\simeq S^1$ and
$BM\simeq S^2$. Let $B$ have two objects $a,b$ and two
nonidentity arrows $f,g:a\rightrightarrows b$. Define $u:Y\to B$
by minima to $a$, maxima to $b$, $r_0\to q_1$ to $g$, and the
other three arrows to $f$. The cycle
\[
 \gamma=(r_0q_0)-(r_1q_0)+(r_1q_1)-(r_0q_1)
\]
maps to $f-g$. These graphs are circles, so $Bu$ has degree one
and $u\in W_\infty$.

Because $M$ is a poset, a functor $B\to M$ sends $f$ and $g$
to the same arrow, and the same equation identifies natural
transformations. Thus $M^B\cong M^{[1]}$. Evaluation at the
source has a constant-arrow section, with a natural transformation
from section-after-evaluation to the identity. Its nerve is
therefore homotopy equivalent to $BM$. Since every arrow is
marked, ordinary localization is groupoid completion, so
$H_{s,s}(B)=\pi_1 B(M^B,s_B)=0$.

In $M^Y$ let $j:Y\hookrightarrow M$ be inclusion and choose
$y_0=r_0$. The loop
\[
 s_Y\longleftarrow j\longrightarrow n_Y
       \longleftarrow (y_0)_Y\longrightarrow s_Y
                                                               \tag{C.25.1}
\]
is well typed: the two last arrows use $y_0\to n,s$. Evaluation
on $BY$ turns its first two legs into the two cone homotopies on
the equatorial circle. Its last two legs are constant in the
$BY$ variable. Therefore the image of $[\gamma]\times[S^1]$
is the cycle $n*\gamma-s*\gamma$. This can be checked on
each edge prism: its cone triangle has the indicated orientation,
and the fixed return leg has no two-dimensional image. The eight
triangles have coefficients $\pm1$; their boundary is zero and
the poset has no three-simplices. Pairing with its fundamental
cohomology class gives a homomorphism
$\pi_1B(M^Y,s_Y)\to\mathbb Z$ sending (C.25.1) to one.
Concatenation adds the edge-prism chains, proving this is a
homomorphism. The loop has infinite order, whereas its prospective
source under $H_{s,s}(u)$ is the trivial group. Hence restriction
is not surjective. The independent program verifies the relation
tables, both nerve complexes, the eight triangle coefficients and
boundary, and counts $198$ functors $Y\to M$ and $18$ functors
$B\to M$.

C.7 gives $B\operatorname{Ch}(M;s,s)\simeq\Omega S^2$ in this
same example. Every category has a Thomason-fibrant replacement
by a specified weak equivalence. Thus the reading ``isomorphic in
$\mathrm{Hot}$ to a Thomason-fibrant category'' holds here while
invariance fails. The first sentence of this claim follows from
these exact comparisons, rather than from a terminological
distinction.
\end{proof}

\begin{claim}{C.26}{Refuted, with two complete replacements}{IX ms.~22--23: the finite example and the loop test}
The two-object category of pointed sets on the one- and two-element
pointed sets, marked only at identities, satisfies both
representabilities in test (a), global invariance in test (b), and
both pointed mapping-space formulas in test (c). Its localization
is not equivalent to a point. Thus the finite-category doubt at
ms.~23 is false under the hypotheses printed there.

Conversely there is a countable small relative category satisfying
both formulas in (c) for all objects while global invariance in
(b) fails. Thus the final assertion on ms.~23 that (b) is
necessary for (c) is also false. The exact positive statement is
that genuine homotopy pullbacks $\Omega b=0\times_b0$ and
homotopy pushouts $\Sigma a=0\amalg_a0$ in the localized
$\infty$-category give these formulas with all their loop
coherences, by the mapping-space universal properties.
\end{claim}
\begin{proof}
The finite category $E$ has arrows
\[
 1_0,\quad u:0\to x,\quad v:x\to0,\quad1_x,
 \quad z=uv:x\to x,
 \qquad vu=1_0,\quad z^2=z\ne1_x.
\]
All composites are the composites of these actual pointed-set
maps, so this is a complete five-arrow category. The object zero
is both initial and terminal. Localization at identities preserves
$E$ and every $E^X$. Naturality of a transformation between
constant diagrams says exactly that its selected arrow is constant
on each zigzag component of $X$; hence
\[
 H_{a,b}(X)=\operatorname{Hom}_E(a,b)^{\pi_0X}.
\]
This proves global invariance, including the empty test. Moreover
$\operatorname{Fun}(NX,NE)=N\operatorname{Fun}(X,E)$: in
degree $r$, both are the functors $X\times[r]\to E$, with the
same simplicial operations. Thus the comparison of diagram
localizations is an actual isomorphism here.

All localization mapping spaces are discrete. For the actual
mixed-block path systems, delete each negative edge, which is
forced to be an identity. This is a category isomorphism to the
positive path category of length equal to the number of positive
edges. Its inverse repeats the deleted vertices; the naturality
square at an identity forces equality of the repeated vertical
components. These isomorphisms commute with the left, right and
symmetric padding maps, by counting the positive edges in each
block. C.5 therefore gives categorical homotopy equivalences of
their colimits with the discrete Hom sets. This verifies even the
strong local reading of the first version of test (b).
The based groups $\pi_1$ are trivial, represented in both variables
by the zero object. Set $\Sigma a=0=\Omega a$. Both sides of
each formula in (c) are contractible, with their unique group
structures. The distinct arrows $1_x,z$ survive localization,
so $E$ is not equivalent to the terminal category. The program
checks all $34$ composable triples and the zero-object property.

For the other direction take the small category $\mathcal M$ of
bounded chain complexes of finite-dimensional rational vector
spaces, represented by finite rational matrices, and mark the
quasi-isomorphisms. There are countably many objects and arrows.
Use degreewise injections as cofibrations and degreewise
surjections as fibrations. Here is a check of the precise
model-category hypotheses of \cite[\S1.4(iv)]{CpathsDKFunction},
whose axiom M0 asks only for finite direct and inverse limits.
Those limits exist in $\mathcal M$. Every acyclic complex of
vector spaces is contractible, by choosing complements to its
cycles. In a lifting problem, choose a lift of graded vector
spaces first. Its differential defect is a cycle in
$\operatorname{Hom}_*(B/A,\ker p)$. If the injection or the
surjection is a quasi-isomorphism, the respective quotient or
kernel is contractible, so this Hom complex is contractible.
Correcting the graded lift by a boundary gives a chain-map lift.

For a map $f:A\to B$, put $C(A)=\operatorname{Cone}(1_A)$,
with its chain injection $j:A\to C(A)$. The factorization
\[
 A\xrightarrow{(f,j)}B\oplus C(A)\longrightarrow B
\]
is a cofibration followed by a trivial fibration. The natural
degreewise surjection $q:C(B)[-1]\to B$ gives the other
factorization
\[
 A\longrightarrow A\oplus C(B)[-1]
                    \xrightarrow{(f,q)}B,
\]
a trivial cofibration followed by a fibration. Finite limits
preserve the relevant kernels, finite colimits preserve the
relevant cokernels, and exactness of vector spaces verifies their
base-change axioms. Two-out-of-three follows from homology.
This checks every axiom used by the cited function-complex theorem.
All objects are both cofibrant and fibrant.

The cosimplicial object $A\otimes N\mathbb Q\Delta^\bullet$
is a cosimplicial resolution of $A$: its values augment by
quasi-isomorphisms, and its latching maps are the degreewise
injections associated with the boundaries of simplices.
Proposition~4.4 of that source identifies the specified function
complex with the simplicial abelian group
\[
 r\longmapsto
 \operatorname{Hom}_{\mathrm{Ch}}
       (A\otimes N\mathbb Q\Delta^r,B).
\]
Its normalized chain complex is
$\tau_{\geq0}\operatorname{Hom}_*(A,B)$, using cycles in degree
zero: by the tensor--Hom adjunction an $r$-simplex is exactly a
chain map $N\mathbb Q\Delta^r\to\operatorname{Hom}_*(A,B)$,
and its normalized nondegenerate value is the image of the top
simplex, with differential its alternating face sum. This is the
explicit Dold--Kan calculation. Therefore
\[
 \pi_i\operatorname{Map}_{L(\mathcal M,W)}(A,B)
       =H_i\operatorname{Hom}_*(A,B),\qquad i\geq0.
\]
The calculation is natural in both complexes and under their
composition. Shifts $\Sigma A=A[1]$ and $\Omega B=B[-1]$
shift this Hom complex down once; its good nonnegative truncation
models the based loop simplicial abelian group. Equivalently the
cone factorizations above compute $0\times_B^h0=B[-1]$ and
$0\amalg_A^h0=A[1]$. The mapping-space universal properties
then give both formulas in (c), with the iterated pullback
coherences of loop multiplication.

At $A=\mathbb Q[0]$, $B=\mathbb Q[2]$ the mapping space is
$K(\mathbb Q,2)$. Global invariance would contradict B.1 and
C.20: the invariant ind-path functor depends only on $\Pi_1$;
its descended value is $\operatorname{colim}_n[BX,BK_n]$ by
(B.1.1). Test a finite face-poset model of $S^2$ and the point.
Compactness of the finite sphere and the specified C.3 comparison
identify these colimits with $[S^2,K(\mathbb Q,2)]$ and
$[*,K(\mathbb Q,2)]$, respectively. They are $\mathbb Q$ and
a singleton, although both test categories have terminal
fundamental groupoid. This is the contradiction. It refutes the
necessity assertion without imposing a comparison hypothesis on
the strict diagram localizations.
\end{proof}
